\scanpage{166}
\mspage{1 - 36 -}
\margin{n° 413}

\uline{avoir de confusion sur} $\varphi$), \uline{l'image par} $\varphi$ \uline{de la distribution ponctuelle} $t_1\otimes\ldots\otimes t_n$ \uline{sur} $X$.

Il résulte immédiatement de ce que nous avons rappelé plus haut que l'application $(t_i)\longrightarrow t_1*\ldots*t_n$ est un morphisme multilinéaire de fibrés vectoriels. Elle s'étend par multilinéarité en une application notée de la même manière de $\prod\mathcal T^{(s_i)}(X_i)$ dans $\mathcal T^{(s)}(X)$.

PROPOSITION 19.- \uline{Soient} $X_{ij}$, $Y_i$ \uline{et} $Z$ \uline{des variétés de classe} $C^r$ (\uline{pour} $1\leq i\leq m$ \uline{et} $1\leq j\leq n_i$), $\varphi_i$ \uline{un morphisme du produit} $X_i=\prod_j X_{ij}$ \uline{dans} $Y_i$, $\varphi$ \uline{le produit des} $\varphi_i$ \uline{et} $\psi$ \uline{un morphisme du produit} $Y$ \uline{des} $Y_i$ \uline{dans} $Z$. \uline{Soient} $t_{ij}$ \uline{des distributions ponctuelles d'ordre} $s_{ij}$ \uline{sur} $X_{ij}$ \uline{et supposons que} $\sum s_{ij}\leq r$. \uline{On a alors} :
\[
(21)\qquad
\underset{\psi\circ\varphi}{*}(t_{ij})_{(1\leq i\leq m,\,\unsure{\leq}1\leq j\leq n_i)}
=\underset{\psi}{*}\bigl(\underset{\varphi_i}{*}(t_{ij})_{(1\leq j\leq n_i)}\bigr)_{(1\leq i\leq m)}.
\]
Cela résulte aussitôt des formules (18) et (20).

COROLLAIRE.- \uline{Soient} $X_i$, $X'_i$ $(1\leq i\leq n)$, $Y$, $Y'$ \uline{des variétés de classe} $C^r$ ; \uline{soient} $\varphi$, $\varphi'$ \uline{des applications de classe} $C^r$ \uline{du produit} $X$ \uline{des} $X_i$ \uline{dans} $Y$ \uline{et du produit} $X'$ \uline{des} $X'_i$ \uline{dans} $Y'$ \uline{respectivement} ; \uline{soient} $f_i$ \uline{des applications de classe} $C^r$ \uline{de} $X_i$ \uline{dans} $X'_i$ $(1\leq i\leq n)$ \uline{et} $g$ \uline{une application de classe} $C^r$ \uline{de} $Y$ \uline{dans} $Y'$. \uline{Soit} $f$ \uline{le produit des} $f_i$ \uline{et soient} $t_i$ \uline{des distributions ponctuelles sur} $X_i$, \uline{d'ordre} $\leq s_i$.

\uline{On suppose que} $\varphi'\circ f=g\circ\varphi$ \uline{et que} $s=\sum s_i$ \uline{est inférieur à} $r$. \uline{On a alors} :
\[
(22)\qquad
\underset{\varphi'}{*}(f_{i*}(t_i))=g\bigl(\underset{\varphi}{*}t_i\bigr)=\underset{g\circ\varphi}{*}(t_i).
\]
C'est un cas particulier de (21).
\margin{165}
